%% theory_directions.tex
%% Theory note: orbit filters, symbol structure, and how few directions an
%% operator really needs.
%%
%% Companion to jsc_paper_main.tex.  Self-contained; every statement is proved.
%% Build sequence: pdflatex, pdflatex.

\documentclass[final]{siamltex}

\usepackage{amsmath,amssymb}
\usepackage{mathrsfs}
\usepackage{mathtools}
\usepackage{bm}
\usepackage{booktabs}
\usepackage{hyperref}
\usepackage{cleveref}
\usepackage{aliascnt}
\usepackage{enumitem}

\hypersetup{hidelinks}

\DeclareMathOperator{\spanop}{span}
\DeclareMathOperator{\Cat}{Cat}
\DeclareMathOperator{\Stab}{Stab}

\newaliascnt{remark}{theorem}
\newtheorem{remark}[remark]{Remark}
\aliascntresetthe{remark}
\crefname{remark}{Remark}{Remarks}
\Crefname{remark}{Remark}{Remarks}

\newaliascnt{example}{theorem}
\newtheorem{example}[example]{Example}
\aliascntresetthe{example}
\crefname{example}{Example}{Examples}
\Crefname{example}{Example}{Examples}

\newcommand{\R}{\mathbb{R}}
\newcommand{\C}{\mathbb{C}}
\newcommand{\N}{\mathbb{N}}
\newcommand{\Z}{\mathbb{Z}}
\newcommand{\K}{\mathbb{K}}
\newcommand{\bx}{\bm{x}}
\newcommand{\by}{\bm{y}}
\newcommand{\bz}{\bm{z}}
\newcommand{\bv}{\bm{v}}
\newcommand{\bw}{\bm{w}}
\newcommand{\Tp}{T_{p}}
\newcommand{\abs}[1]{\left|#1\right|}
\newcommand{\norm}[1]{\left\|#1\right\|}

\emergencystretch=1.5em

\begin{document}

\title{How Few Directions Does an Operator Need?\\
Orbit Filters, Symbol Structure, and Shared Schedules}

\author{Theory note accompanying\\
\emph{Rank-Optimal Directional Taylor Jets for High-Order Mixed Derivatives}}

\maketitle

\begin{abstract}
  The directional schedule of the main paper is built from roots of unity, and
  it is minimal for a single monomial derivative.  This note asks what that
  construction is really doing and how far it goes.  Three answers follow.
  First, the roots-of-unity grid is not specific to monomials at all: for an
  arbitrary homogeneous constant-coefficient operator it produces an exact
  schedule whose coefficients are a discrete Fourier transform of the symbol,
  with no linear solve and no dimension restriction.  Second, that grid is a
  complete intersection, and complete intersections are provably far from
  minimal once three or more coordinates are active; for a generic symbol the
  grid costs about $N!$ times the minimum in $N$ active variables.  Third, both
  the roots-of-unity schedule and its correct generalization are instances of a
  single principle: a character-weighted orbit of a finite group reproduces any
  symbol lying in a one-dimensional isotypic component.  Choosing an abelian
  group recovers the Fourier grid; choosing a nonabelian group turns the design
  problem for powers of the Laplacian into the classical problem of cubature on
  the sphere, whose solutions are real, positively weighted, and in the smallest
  cases provably minimal.  We also record that scheduling several operators, or
  a whole derivative tensor, on one shared direction set is exactly optimal at
  length equal to the dimension of the span, which quantifies how much
  term-by-term scheduling wastes.
\end{abstract}

%% ====================================================================
\section{Scope and notation}\label{sec_scope}

Throughout, $\K \in \{\R,\C\}$ and $\bz = (z_{0},\ldots,z_{n})$ are coordinates
on the active subspace of the problem, so that $N := n+1$ is the number of
active variables.  We write $\K[\bz]_{p}$ for the homogeneous polynomials of
degree $p$, whose dimension is
\begin{equation}\label{eq_dim}
  D_{p,N} \;:=\; \dim \K[\bz]_{p} \;=\; \binom{p+N-1}{N-1}.
\end{equation}
A homogeneous constant-coefficient operator of order $p$ and its symbol are
\begin{equation}\label{eq_symbol}
  \mathcal{P} \;=\; \sum_{\abs{\beta}=p} b_{\beta}\,\partial^{\beta},
  \qquad
  \sigma(\bz) \;=\; \sum_{\abs{\beta}=p} b_{\beta}\,\bz^{\beta} \in \K[\bz]_{p}.
\end{equation}
For a direction $\bv$ we write $\bv \cdot \bz = \sum_{j} v_{j} z_{j}$, and
$\Tp(\bx;\bv)$ for the order-$p$ coefficient of $\tau \mapsto u(\bx + \tau\bv)$,
as in the main paper.  A \emph{directional schedule} for $\mathcal{P}$ of length
$R$ is a family $(c_{r},\bv_{r})_{r=1}^{R}$ with
$\mathcal{P}u(\bx) = \sum_{r} c_{r}\,\Tp(\bx;\bv_{r})$ for all admissible $u$.
We take the following bridge as given; it is Proposition~3.6 of the main paper.

\begin{proposition}[Bridge]\label{prop_bridge}
  A family $(c_{r},\bv_{r})_{r=1}^{R}$ is a directional schedule for
  $\mathcal{P}$ if and only if
  \begin{equation}\label{eq_bridge}
    p!\,\sigma(\bz) \;=\; \sum_{r=1}^{R} c_{r}\,(\bv_{r}\cdot\bz)^{p}
    \qquad\text{in }\K[\bz]_{p},
  \end{equation}
  and the least length of a schedule over $\K$ is the Waring rank
  $R_{\K}(\sigma)$.
\end{proposition}

Everything below is therefore a statement about writing one fixed form $\sigma$
as a short combination of $p$th powers of linear forms.  Two standard devices
are used.  Let $\K[\by] = \K[y_{0},\ldots,y_{n}]$ act on $\K[\bz]$ by
$y_{j} \mapsto \partial/\partial z_{j}$.  The \emph{apolar ideal} of $\sigma$ is
\begin{equation}\label{eq_apolar}
  \sigma^{\perp} \;:=\; \{\Theta \in \K[\by] \;:\; \Theta(\partial)\,\sigma = 0\},
\end{equation}
a homogeneous ideal, and the \emph{catalecticant} in degree $q$ is the linear map
\begin{equation}\label{eq_cat}
  \Cat_{q}(\sigma) \colon \K[\by]_{q} \to \K[\bz]_{p-q},
  \qquad
  \Theta \mapsto \Theta(\partial)\sigma,
  \qquad
  \ker \Cat_{q}(\sigma) = (\sigma^{\perp})_{q}.
\end{equation}
We use the apolarity lemma in the form: for a finite reduced
$X \subset \mathbb{P}(\K^{N})$ with homogeneous vanishing ideal $I_{X}$,
\begin{equation}\label{eq_apolarity}
  \sigma \in \spanop\{(\bv\cdot\bz)^{p} : [\bv] \in X\}
  \iff
  I_{X} \subseteq \sigma^{\perp}.
\end{equation}
Finally, $U_{m} := \{\zeta \in \C : \zeta^{m}=1\}$, and we use repeatedly that
\begin{equation}\label{eq_filter}
  \sum_{\zeta \in U_{m}} \zeta^{k} \;=\; m \cdot \mathbf{1}_{m \mid k},
  \qquad k \in \Z .
\end{equation}

%% ====================================================================
\section{The Fourier grid is not about monomials}\label{sec_grid}

The schedule of the main paper attaches roots of unity to a monomial
$\bz^{a}$.  The next theorem shows that the same point set serves an arbitrary
symbol, with coefficients given in closed form by a discrete Fourier transform.
No linear system is solved, and no restriction is placed on $N$.

\begin{theorem}[Universal Fourier schedule]\label{thm_grid}
  Let $\sigma \in \K[\bz]_{p}$ be nonzero with coefficients $b_{\beta}$ as in
  \eqref{eq_symbol}.  Put
  \begin{equation}\label{eq_degrees}
    d_{j} := \deg_{z_{j}} \sigma,
    \qquad
    m_{j} := d_{j}+1,
    \qquad j = 0,\ldots,n,
  \end{equation}
  and relabel the coordinates so that $m_{0} = \min_{j} m_{j}$.  Set
  $M := \prod_{j=1}^{n} m_{j}$.  For
  $\zeta = (\zeta_{1},\ldots,\zeta_{n}) \in U_{m_{1}} \times \cdots \times U_{m_{n}}$
  define
  \begin{equation}\label{eq_grid}
    \bv_{\zeta} \;:=\; e_{0} + \sum_{j=1}^{n} \zeta_{j}\,e_{j},
    \qquad
    c_{\zeta} \;:=\; \frac{1}{M}\sum_{\abs{\gamma}=p} \gamma!\,b_{\gamma}
      \prod_{j=1}^{n} \zeta_{j}^{-\gamma_{j}} .
  \end{equation}
  Then $(c_{\zeta},\bv_{\zeta})_{\zeta}$ is a directional schedule for
  $\mathcal{P}$ of length $M$.
\end{theorem}

\begin{proof}
  By \cref{prop_bridge} it suffices to verify \eqref{eq_bridge}.  Since
  $\bv_{\zeta} \cdot \bz = z_{0} + \sum_{j \ge 1}\zeta_{j}z_{j}$, the
  multinomial theorem gives
  \begin{equation*}
    \sum_{\zeta} c_{\zeta}(\bv_{\zeta}\cdot\bz)^{p}
    = \sum_{\abs{\beta}=p} \frac{p!}{\beta!}\, \bz^{\beta}
      \underbrace{\sum_{\zeta} c_{\zeta} \prod_{j=1}^{n}\zeta_{j}^{\beta_{j}}}_{=:\,S_{\beta}},
  \end{equation*}
  so \eqref{eq_bridge} is equivalent to $S_{\beta} = \beta!\,b_{\beta}$ for every
  $\beta$ with $\abs{\beta}=p$.  Substituting \eqref{eq_grid} and applying
  \eqref{eq_filter} in each factor,
  \begin{equation}\label{eq_Sbeta}
    S_{\beta}
    = \frac{1}{M}\sum_{\abs{\gamma}=p}\gamma!\,b_{\gamma}
      \prod_{j=1}^{n}\Bigl(\sum_{\zeta_{j}\in U_{m_{j}}}\zeta_{j}^{\beta_{j}-\gamma_{j}}\Bigr)
    = \sum_{\abs{\gamma}=p}\gamma!\,b_{\gamma}
      \prod_{j=1}^{n}\mathbf{1}_{m_{j}\,\mid\,\beta_{j}-\gamma_{j}} .
  \end{equation}

  We claim that for $\gamma$ with $b_{\gamma}\neq 0$ the product of indicators in
  \eqref{eq_Sbeta} equals $1$ precisely when $\gamma = \beta$.  Sufficiency is
  clear.  For necessity, write $\beta_{j}-\gamma_{j} = k_{j}m_{j}$ with
  $k_{j}\in\Z$ for $j = 1,\ldots,n$.  Because $b_{\gamma}\neq 0$, the definition
  \eqref{eq_degrees} forces $\gamma_{j} \le d_{j} = m_{j}-1$ for \emph{every}
  $j$, including $j=0$.  If $k_{j} \le -1$ for some $j \ge 1$, then
  $\gamma_{j} \ge \beta_{j}+m_{j} \ge m_{j} > d_{j}$, a contradiction; hence
  $k_{j}\ge 0$ for all $j \ge 1$.  Summing $\abs{\beta} = \abs{\gamma} = p$ gives
  $\sum_{j\ge 0}(\beta_{j}-\gamma_{j}) = 0$, that is,
  \begin{equation}\label{eq_alias}
    \sum_{j=1}^{n} k_{j}m_{j} \;=\; \gamma_{0}-\beta_{0} \;\le\; \gamma_{0} \;\le\; d_{0} \;=\; m_{0}-1 .
  \end{equation}
  If some $k_{j} \ge 1$ then the left side of \eqref{eq_alias} is at least
  $m_{j} \ge m_{0}$, contradicting the bound $m_{0}-1$.  Therefore $k_{j}=0$ for
  all $j \ge 1$, so $\beta_{j}=\gamma_{j}$ for $j\ge1$, and then
  $\beta_{0}=\gamma_{0}$ because $\abs{\beta}=\abs{\gamma}$.  This proves the
  claim.

  Consequently the sum \eqref{eq_Sbeta} retains at most the single term
  $\gamma=\beta$.  If $b_{\beta}\neq0$ it retains exactly that term and
  $S_{\beta} = \beta!\,b_{\beta}$.  If $b_{\beta}=0$, no $\gamma$ with
  $b_{\gamma}\neq0$ survives, since surviving would force $\gamma=\beta$ and
  hence $b_{\beta}\neq0$; thus $S_{\beta}=0=\beta!\,b_{\beta}$.
\end{proof}

Two features of \cref{thm_grid} deserve comment.

\begin{remark}[The base coordinate must be one of least degree]\label{rmk_base}
  The relabelling $m_{0}=\min_{j}m_{j}$ is not cosmetic; it is exactly what
  \eqref{eq_alias} consumes.  Take $\sigma = z_{0}^{4} + z_{0}z_{1}^{3}$, so
  $p=4$, $d_{0}=4$, $d_{1}=3$ and the least degree occurs at $j=1$.  Building
  the grid on the base $j=0$ instead uses $\zeta \in U_{4}$ and the pair
  $\beta=(0,4)$, $\gamma=(4,0)$ satisfies $4 \mid \beta_{1}-\gamma_{1}$ with
  $\gamma \neq \beta$.  The aliased term survives and the construction returns
  $\sigma + z_{1}^{4}$ rather than $\sigma$.  Choosing the base correctly makes
  every alias impossible at once.
\end{remark}

\begin{remark}[Recovering the monomial case]\label{rmk_monomial}
  For $\sigma = \bz^{a}$ with $a$ the active exponent vector, only
  $\gamma = a$ contributes to \eqref{eq_grid}, so
  $c_{\zeta} = (a!/M)\prod_{j}\zeta_{j}^{-a_{j}}$ and $M = \prod_{j\ge1}(a_{j}+1)$.
  This is Theorem~3.4 of the main paper, up to replacing $\zeta_{j}$ by
  $\zeta_{j}^{-1}$, which permutes the grid.  So the monomial schedule is the
  one-term instance of \cref{thm_grid}.
\end{remark}

The practical reading of \cref{thm_grid} is that the grid is a
\emph{universal} device: it needs nothing from $\sigma$ beyond its coordinate
degrees, it never solves a linear system, and its length
$M = \prod_{j\ge1}m_{j}$ is available in advance.  The rest of this note is
about the price of that universality.

%% ====================================================================
\section{What grid structure costs}\label{sec_cost}

The grid of \cref{thm_grid} is the zero set of the $n$ binomials
$y_{j}^{m_{j}} - y_{0}^{m_{j}}$, $j = 1,\ldots,n$, hence a complete intersection
in $\mathbb{P}^{n}$.  That single structural fact bounds how good it can ever be.

\begin{definition}[Initial degree]\label{def_initial}
  For $0 \neq \sigma \in \K[\bz]_{p}$ set
  $\delta(\sigma) := \min\{q \ge 1 : (\sigma^{\perp})_{q} \neq 0\}$.
\end{definition}

\begin{proposition}[Complete intersections are large]\label{prop_ci}
  Let $X \subset \mathbb{P}^{n}$ be a finite apolar scheme for $\sigma$, that is
  $I_{X} \subseteq \sigma^{\perp}$, and suppose $X$ is a complete intersection of
  $n$ forms of degrees $e_{1},\ldots,e_{n}$.  Then
  \begin{equation}\label{eq_ci}
    \abs{X} \;=\; \prod_{i=1}^{n} e_{i} \;\ge\; \delta(\sigma)^{\,n}.
  \end{equation}
\end{proposition}

\begin{proof}
  The $n$ defining forms lie in $I_{X} \subseteq \sigma^{\perp}$, so each has
  degree at least $\delta(\sigma)$ by \cref{def_initial}.  Since $X$ is a
  complete intersection of these forms, B\'ezout's theorem gives
  $\deg X = \prod_{i} e_{i}$, and $X$ is finite and reduced so
  $\abs{X} = \deg X$.
\end{proof}

It remains to locate $\delta(\sigma)$ for a symbol with no special structure.

\begin{lemma}[Initial degree of a generic symbol]\label{lem_generic_delta}
  For $\sigma$ generic in $\K[\bz]_{p}$ one has
  $\delta(\sigma) = \lfloor p/2 \rfloor + 1$.
\end{lemma}

\begin{proof}
  By \eqref{eq_cat}, $(\sigma^{\perp})_{q} \neq 0$ if and only if
  $\Cat_{q}(\sigma)$ has nontrivial kernel.  The rank of $\Cat_{q}$ is a lower
  semicontinuous function of $\sigma$, so it suffices to exhibit one $\sigma$
  for which $\Cat_{q}$ has the maximal possible rank
  $\min(D_{q,N}, D_{p-q,N})$ for every $q$; genericity then gives maximal rank
  for all $q$ simultaneously.  Take $\sigma = \sum_{i=1}^{s}\ell_{i}^{p}$ with
  $s \ge \max_{q}\min(D_{q,N},D_{p-q,N})$ and $\ell_{i}$ generic linear forms.
  Then $\Cat_{q}(\ell^{p})$ is the rank-one map
  $\Theta \mapsto \frac{p!}{(p-q)!}\Theta(\ell)\,\ell^{p-q}$, so
  $\Cat_{q}(\sigma)$ has image inside $\spanop\{\ell_{i}^{p-q}\}$ and rank equal
  to the rank of the $s \times D_{q,N}$ evaluation matrix
  $(\Theta_{k}(\ell_{i}))$ composed with the independence of
  $\{\ell_{i}^{p-q}\}$.  For generic $\ell_{i}$ both are of maximal rank, because
  the Veronese variety is nondegenerate and therefore no nonzero form of degree
  $q$ vanishes at $D_{q,N}$ generic points, and likewise
  $\{\ell_{i}^{p-q}\}$ spans $\K[\bz]_{p-q}$ once $s \ge D_{p-q,N}$.  Hence
  $\Cat_{q}$ has maximal rank, and its kernel is nonzero exactly when
  $D_{q,N} > D_{p-q,N}$, that is when $q > p-q$.  The least such $q$ is
  $\lfloor p/2\rfloor+1$.
\end{proof}

\begin{corollary}[The price of grid structure]\label{cor_price}
  Fix $N \ge 2$ and let $\sigma \in \K[\bz]_{p}$ be generic.  Every complete
  intersection apolar scheme for $\sigma$ has at least
  $(\lfloor p/2\rfloor+1)^{N-1}$ points, while $R_{\C}(\sigma)$ is asymptotic to
  $D_{p,N}/N$ as $p \to \infty$.  Consequently
  \begin{equation}\label{eq_price}
    \lim_{p\to\infty}
    \frac{(\lfloor p/2\rfloor+1)^{N-1}}{R_{\C}(\sigma)}
    \;=\; \frac{N!}{2^{\,N-1}},
    \qquad\text{and}\qquad
    \lim_{p\to\infty}\frac{M}{R_{\C}(\sigma)} \;=\; N!,
  \end{equation}
  where $M = (p+1)^{N-1}$ is the length of the grid of \cref{thm_grid} for a
  generic symbol.
\end{corollary}

\begin{proof}
  The first bound is \cref{prop_ci} together with \cref{lem_generic_delta}.
  The generic Waring rank is $\lceil D_{p,N}/N \rceil$ outside the finite list of
  Alexander--Hirschowitz exceptions, and
  $D_{p,N} = \binom{p+N-1}{N-1} \sim p^{N-1}/(N-1)!$, so
  $R_{\C}(\sigma) \sim p^{N-1}/N!$.  Dividing
  $(\lfloor p/2\rfloor+1)^{N-1} \sim (p/2)^{N-1}$ by this gives
  $N!/2^{N-1}$.  A generic $\sigma$ has $d_{j}=p$ for every $j$, so
  $M = (p+1)^{N-1} \sim p^{N-1}$, and dividing gives $N!$.
\end{proof}

The two limits in \eqref{eq_price} separate two different losses.  The factor
$N!/2^{N-1}$ is the intrinsic cost of insisting on complete intersection
structure, and the remaining factor $2^{N-1}$ is the cost of building the grid
from the coordinate degrees rather than from $\sigma^{\perp}$.  Both vanish when
$N = 2$: there $N!/2^{N-1} = 1$, which is the algebraic reason binary forms are
easy, since every finite subscheme of $\mathbb{P}^{1}$ is cut out by a single
form and the complete intersection hypothesis is vacuous.  For $N=3$ the
intrinsic factor is $3/2$, for $N=4$ it is $3$, for $N=5$ it is $15/2$, and it
grows superexponentially.

\begin{remark}[When the grid is exactly minimal]\label{rmk_tight}
  Suppose $\sigma = \ell_{0}^{a_{0}}\cdots\ell_{n}^{a_{n}}$ for linearly
  independent linear forms $\ell_{j}$, that is, $\sigma$ is projectively
  equivalent to a monomial.  Waring rank is invariant under invertible linear
  changes of variables, so running \cref{thm_grid} in the coordinate system dual
  to $(\ell_{j})$ yields a schedule of length $\prod_{j\ge1}(a_{j}+1)$, which is
  the rank of that monomial and hence minimal.  Corollary~3.7 of the main paper
  is the case $\sigma = (z_{1}^{2}+z_{2}^{2})^{m} = (z_{1}+\mathrm{i}z_{2})^{m}(z_{1}-\mathrm{i}z_{2})^{m}$
  in two variables.  The hypothesis is restrictive: forms equivalent to
  monomials constitute a subvariety of dimension at most $N^{2}$ inside the
  $D_{p,N}$-dimensional space of all symbols.
\end{remark}

%% ====================================================================
\section{Orbit filters}\label{sec_orbit}

Identity \eqref{eq_filter} is character orthogonality for the cyclic group
$U_{m}$.  Read that way, the roots-of-unity construction is one instance of a
principle that does not require the group to be abelian, and the nonabelian
instances are the interesting ones.  Throughout this section $G$ is a finite
subgroup of the orthogonal group $O(N)$, acting on $\K[\bz]_{p}$ by
$(g \cdot F)(\bz) := F(g^{-1}\bz)$.  Orthogonality gives the compatibility we
need,
\begin{equation}\label{eq_compat}
  (g\bv)\cdot\bz \;=\; \bv \cdot (g^{-1}\bz),
\end{equation}
and we write $\pi_{G} := \abs{G}^{-1}\sum_{g\in G} g$ for the Reynolds operator,
a projection of $\K[\bz]_{p}$ onto the invariants $\K[\bz]_{p}^{G}$.

\begin{definition}[Orbit average]\label{def_orbit_avg}
  For $\bv \in \K^{N}$ and a linear character $\chi \colon G \to \K^{\times}$ put
  \begin{equation}\label{eq_orbit_avg}
    A^{\chi}_{\bv}(\bz) \;:=\; \frac{1}{\abs{G}}\sum_{g \in G}
    \overline{\chi(g)}\,\bigl((g\bv)\cdot\bz\bigr)^{p},
  \end{equation}
  and write $A_{\bv} := A^{\mathbf{1}}_{\bv}$ for the trivial character.  The
  \emph{isotypic component} of $\chi$ is
  $\K[\bz]_{p}^{(\chi)} := \{F : g\cdot F = \chi(g)F \text{ for all } g\}$; for
  $\chi = \mathbf{1}$ it is $\K[\bz]_{p}^{G}$.
\end{definition}

\begin{lemma}[Orbit averages are isotypic]\label{lem_isotypic}
  $A^{\chi}_{\bv} \in \K[\bz]_{p}^{(\chi)}$ for every $\bv$ and every linear
  character $\chi$.  Moreover $A^{\chi}_{\bv}$ is the image of
  $(\bv\cdot\bz)^{p}$ under the projection
  $\pi^{\chi}_{G} := \abs{G}^{-1}\sum_{g}\overline{\chi(g)}\,g$ onto
  $\K[\bz]_{p}^{(\chi)}$, and consequently
  $\{A^{\chi}_{\bv} : \bv \in \K^{N}\}$ spans $\K[\bz]_{p}^{(\chi)}$.
\end{lemma}

\begin{proof}
  Using \eqref{eq_compat} and then reindexing $g \mapsto h^{-1}g$,
  \begin{equation*}
    (h\cdot A^{\chi}_{\bv})(\bz)
    = \frac{1}{\abs{G}}\sum_{g}\overline{\chi(g)}\bigl((g\bv)\cdot h^{-1}\bz\bigr)^{p}
    = \frac{1}{\abs{G}}\sum_{g}\overline{\chi(g)}\bigl((hg\bv)\cdot \bz\bigr)^{p}
    = \chi(h)\,A^{\chi}_{\bv}(\bz),
  \end{equation*}
  where the last step substitutes $g \mapsto h^{-1}g$ and uses
  $\overline{\chi(h^{-1}g)} = \chi(h)\overline{\chi(g)}$ for a character of
  modulus one.  The same computation with \eqref{eq_compat} shows
  $\pi^{\chi}_{G}\bigl((\bv\cdot\bz)^{p}\bigr) = A^{\chi}_{\bv}$.  Since powers
  of linear forms span $\K[\bz]_{p}$ and $\pi^{\chi}_{G}$ is a surjection onto
  $\K[\bz]_{p}^{(\chi)}$, the images $A^{\chi}_{\bv}$ span that component.
\end{proof}

\begin{theorem}[Orbit filter]\label{thm_orbit}
  Let $\bv \in \K^{N}$ with orbit $X := G\bv$, and let $\chi$ be a linear
  character of $G$.
  \begin{enumerate}[label=\textup{(\roman*)}]
    \item If $\sigma \in \K[\bz]_{p}^{(\chi)}$ and $\sigma$ is a scalar multiple
      of $A^{\chi}_{\bv} \neq 0$, then $\mathcal{P}$ admits a directional
      schedule supported on $X$, with coefficients
      $c_{g\bv} = \lambda\,\overline{\chi(g)}$ constant along $\chi$ up to the
      character.  Its length is
      $\abs{X} = \abs{G}/\abs{\Stab(\bv)}$, and $\abs{X}/2$ if $p$ is even and
      $-I \in G$.
    \item In particular, if $\dim_{\K}\K[\bz]_{p}^{(\chi)} = 1$ then the
      hypothesis of \textup{(i)} holds for every $\sigma$ in that component and
      every $\bv$ with $A^{\chi}_{\bv} \neq 0$, with no equation to solve.
    \item If $\sigma = \sum_{\chi} \sigma^{(\chi)}$ is the decomposition of a
      general $\sigma$ into components over the linear characters, and each
      $\sigma^{(\chi)}$ satisfies the hypothesis of \textup{(i)} for the same
      $\bv$ with multiplier $\lambda_{\chi}$, then $\mathcal{P}$ admits a
      schedule supported on the \emph{same} set $X$, with
      $c_{g\bv} = \sum_{\chi}\lambda_{\chi}\overline{\chi(g)}$.
  \end{enumerate}
\end{theorem}

\begin{proof}
  (i) Write $\sigma = \mu\,A^{\chi}_{\bv}$ with $\mu \in \K^{\times}$.  By
  \eqref{eq_orbit_avg},
  $p!\,\sigma = \frac{p!\mu}{\abs{G}}\sum_{g}\overline{\chi(g)}((g\bv)\cdot\bz)^{p}$,
  which is \eqref{eq_bridge} with $\lambda := p!\mu/\abs{G}$; collecting the
  $\abs{\Stab(\bv)}$ group elements that produce the same direction turns the
  sum over $G$ into a sum over $X$.  If $p$ is even and $-I \in G$, then
  $((-\bw)\cdot\bz)^{p} = (\bw\cdot\bz)^{p}$, so antipodal members of $X$ carry
  identical powers and may be merged, halving the count.  (ii) is immediate from
  (i) because a one-dimensional space contains no directions other than
  multiples of a fixed nonzero vector, and $\sigma \neq 0$ forces
  proportionality.  (iii) follows by summing the identities of (i) over $\chi$,
  which changes only the coefficients and not the support.
\end{proof}

\cref{thm_orbit} is the statement that a symmetry group acts as a filter: the
$\binom{p+n}{n}$ moment conditions of the design problem collapse to
$\dim \K[\bz]_{p}^{(\chi)}$ conditions per character, and when that dimension is
one there are no conditions at all.  The two constructions of this note are the
two extremes of the principle.

\begin{example}[The Fourier grid is the abelian case]\label{ex_abelian}
  Let $G := U_{m_{1}} \times \cdots \times U_{m_{n}}$ act diagonally by
  $z_{j} \mapsto \zeta_{j}z_{j}$ for $j \ge 1$ and trivially on $z_{0}$.  Then
  $g\cdot\bz^{\beta} = \bigl(\prod_{j\ge1}\zeta_{j}^{-\beta_{j}}\bigr)\bz^{\beta}$,
  so every monomial spans an isotypic component, and two monomials $\beta$,
  $\beta'$ of degree $p$ share a component if and only if
  $m_{j} \mid \beta_{j}-\beta'_{j}$ for all $j \ge 1$.  When
  $m_{0} = \min_{j}m_{j}$, the argument around \eqref{eq_alias} shows this forces
  $\beta = \beta'$, so \emph{every} component meeting the symbol is
  one-dimensional.  Applying \cref{thm_orbit}(iii) to
  $\bv = e_{0}+e_{1}+\cdots+e_{n}$, whose orbit is exactly the grid of
  \eqref{eq_grid}, recovers \cref{thm_grid}.  The cost of the abelian choice is
  visible here: the group has as many one-dimensional components as the grid has
  points, so the orbit must be as large as the box of monomials it separates,
  which is $\prod_{j\ge1}m_{j}$ rather than the far smaller $D_{p,N}$.
\end{example}

\begin{example}[A nonabelian filter]\label{ex_icosa}
  Let $G \subset O(3)$ be the icosahedral group of order $60$ and let
  $\bv$ be a vertex of a regular icosahedron inscribed in $S^{2}$, whose orbit
  has $12$ points forming $6$ antipodal pairs.  The invariant ring of $G$ acting
  on $\R^{3}$ is generated in degrees $2$, $6$, $10$ and one degree-$15$
  element, so its Molien series gives
  $\dim \R[\bz]_{4}^{G} = 1$, spanned by $(z\cdot z)^{2}$, while
  $\dim\R[\bz]_{4} = 15$.  By \cref{thm_orbit}(ii) the single orbit reproduces
  $(z \cdot z)^{2}$ with equal weights, so $\Delta^{2}$ on $\R^{3}$ has a real
  schedule of length $6$.  \cref{sec_cubature} shows that $6$ is minimal.
\end{example}

%% ====================================================================
\section{Powers of the Laplacian are a cubature problem}\label{sec_cubature}

\cref{ex_icosa} is not a coincidence but an instance of a general equivalence.
Let $q(\bz) := \bz\cdot\bz$ on $\R^{N}$, let $\sigma = q^{m}$ so that
$\mathcal{P} = \Delta^{m}$ and $p = 2m$, and let $\mu$ denote the uniform
probability measure on $S^{N-1}$.

\begin{theorem}[Schedules for $\Delta^{m}$ are spherical cubature rules]\label{thm_cubature}
  Let $\bv_{1},\ldots,\bv_{R} \in S^{N-1}$ and $w_{1},\ldots,w_{R} \in \R$ with
  $\sum_{r} w_{r} = 1$.  The following are equivalent.
  \begin{enumerate}[label=\textup{(\roman*)}]
    \item $\sum_{r=1}^{R} w_{r}(\bv_{r}\cdot\bz)^{2m} = \kappa_{N,m}\,q(\bz)^{m}$
      for all $\bz \in \R^{N}$, where
      $\kappa_{N,m} := (2m-1)!!\,/\,\bigl(N(N+2)\cdots(N+2m-2)\bigr)$;
    \item $\displaystyle\sum_{r=1}^{R} w_{r}f(\bv_{r}) = \int_{S^{N-1}} f\,d\mu$
      for every $f \in \R[\bv]_{2m}$;
    \item the rule $(\bv_{r},w_{r})$ integrates exactly every polynomial of even
      degree at most $2m$ on $S^{N-1}$.
  \end{enumerate}
  In particular a real directional schedule for $\Delta^{m}$ with unit directions
  is the same thing as a cubature rule on $S^{N-1}$ exact in even degrees up to
  $2m$; if in addition the weights are equal and the node set is centrally
  symmetric, it is a spherical $(2m{+}1)$-design.
\end{theorem}

\begin{proof}
  First we identify the constant.  The map
  $\bz \mapsto \int_{S^{N-1}}(\bv\cdot\bz)^{2m}d\mu(\bv)$ is a polynomial of
  degree $2m$, and it is $O(N)$-invariant because $\mu$ is.  The ring of
  $O(N)$-invariant polynomials on $\R^{N}$ is $\R[q]$, so the integral equals
  $\kappa\,q^{m}$ for a constant $\kappa$, which we evaluate at a unit vector
  $\bz = e$: writing $t := \bv \cdot e$ for the first coordinate of a uniform
  point of $S^{N-1}$, $\kappa = \mathbb{E}[t^{2m}]$.  The density of $t$ is
  proportional to $(1-t^{2})^{(N-3)/2}$, and the resulting Beta integral gives
  $\mathbb{E}[t^{2m}] = (2m-1)!!/\bigl(N(N+2)\cdots(N+2m-2)\bigr) = \kappa_{N,m}$.

  (i) $\Leftrightarrow$ (ii).  For fixed $\bz$ the function
  $f_{\bz}(\bv) := (\bv\cdot\bz)^{2m}$ is homogeneous of degree $2m$ in $\bv$,
  and $\int f_{\bz}\,d\mu = \kappa_{N,m}q(\bz)^{m}$ by the previous paragraph.
  Hence (i) states precisely that the rule is exact on every $f_{\bz}$.  Both
  sides of (ii) are linear in $f$, and $\{f_{\bz} : \bz \in \R^{N}\}$ spans
  $\R[\bv]_{2m}$ because powers of linear forms span the space of forms of a
  given degree.  Exactness on a spanning set is exactness on the span, so (i)
  and (ii) are equivalent.

  (ii) $\Leftrightarrow$ (iii).  On $S^{N-1}$ we have $q \equiv 1$, so for
  $k \le m$ and $h \in \R[\bv]_{2k}$ the function $h$ and the form
  $q^{m-k}h \in \R[\bv]_{2m}$ agree at every node and have the same integral;
  thus exactness in degree $2m$ implies exactness in every even degree at most
  $2m$, and the converse is trivial.

  For the final sentence, \cref{prop_bridge} converts (i) into a schedule for
  $\Delta^{m}$ after rescaling the coefficients by
  $(2m)!/\kappa_{N,m}$, and a centrally symmetric equal-weight rule is exact in
  odd degrees automatically, so exactness through degree $2m+1$ holds, which is
  the definition of a spherical $(2m{+}1)$-design.
\end{proof}

\begin{theorem}[Lower bound]\label{thm_lower}
  $R_{\C}(q^{m}) \ge \binom{N+m-1}{m} = D_{m,N}$, and hence every directional
  schedule for $\Delta^{m}$ on $\R^{N}$, real or complex, has at least
  $D_{m,N}$ directions.
\end{theorem}

\begin{proof}
  Consider the symmetric bilinear form on $\R[\by]_{m}$ defined by
  $\langle \Theta,\Psi\rangle := (\Theta\Psi)(\partial)\,q^{m}$, a scalar because
  $\Theta\Psi$ has degree $2m = \deg q^{m}$.  Using the integral representation
  $q^{m} = \kappa_{N,m}^{-1}\int_{S^{N-1}}(\bv\cdot\bz)^{2m}d\mu(\bv)$ from the
  proof of \cref{thm_cubature} and
  $(\Theta\Psi)(\partial)(\bv\cdot\bz)^{2m} = (2m)!\,\Theta(\bv)\Psi(\bv)$,
  \begin{equation*}
    \langle\Theta,\Psi\rangle
    = \frac{(2m)!}{\kappa_{N,m}}\int_{S^{N-1}}\Theta(\bv)\Psi(\bv)\,d\mu(\bv),
  \end{equation*}
  a positive multiple of the $L^{2}(\mu)$ inner product.  It is therefore
  positive definite on $\R[\by]_{m}$, so the middle catalecticant
  $\Cat_{m}(q^{m})$, whose associated bilinear form this is, is nonsingular and
  has rank $D_{m,N}$.

  Now let $q^{m} = \sum_{r=1}^{R}c_{r}(\bv_{r}\cdot\bz)^{2m}$ be any complex
  decomposition.  For $\Theta \in \K[\by]_{m}$,
  $\Theta(\partial)(\bv\cdot\bz)^{2m} = \frac{(2m)!}{m!}\Theta(\bv)(\bv\cdot\bz)^{m}$,
  so the image of $\Cat_{m}(q^{m})$ is contained in
  $\spanop\{(\bv_{r}\cdot\bz)^{m} : r \le R\}$, a space of dimension at most
  $R$.  Comparing with the rank just computed gives $R \ge D_{m,N}$.
\end{proof}

\begin{corollary}[Tight cases]\label{cor_tight}
  For $\Delta^{2}$ on $\R^{3}$ the minimum number of directions is $6$, attained
  over $\R$ by the six icosahedral axes with equal positive weights.  For
  $\Delta^{1}$ on $\R^{N}$ the minimum is $N$, attained by any orthonormal frame.
\end{corollary}

\begin{proof}
  \cref{thm_lower} gives $R \ge D_{2,3} = 6$ and $R \ge D_{1,N} = N$
  respectively.  \cref{ex_icosa} attains the first, and
  $q = \sum_{j}z_{j}^{2}$ attains the second.
\end{proof}

Three consequences are worth stating plainly.  The equivalence removes the
search: for $\Delta^{m}$ one does not look for Waring decompositions but reads
off a cubature rule, of which large tabulated families exist in every dimension.
The nodes are real and the weights are positive, so no complex arithmetic is
needed for polyharmonic operators, in contrast to the roots-of-unity schedule.
And the classical lower bound for centrally symmetric cubature coincides with
the catalecticant bound of \cref{thm_lower}, so the two literatures agree on what
is possible.

\begin{remark}[The bound is not always attained over $\R$]\label{rmk_realgap}
  Attainment is a genuinely finer question than the bound.  A numerical search
  over real nodes, with the weights eliminated by orthogonal projection so that
  only the nodes are optimized, reaches $D_{m,N}$ for $(N,m) = (3,1)$, $(3,2)$
  and $(4,1)$.  For $(N,m) = (3,3)$ and $(4,2)$ it instead stalls at
  $D_{m,N}+1$, namely $11$ nodes against the bound $10$ in both cases, with
  positive weights.  The evidence that the bound is genuinely missed rather than
  merely unfound is that at $R = D_{m,N}$ the residual settles near $10^{-1}$
  over forty restarts, whereas at $R = D_{m,N}+1$ it reaches machine precision
  on the first attempt.  This is consistent with the general fact that real
  Waring rank can exceed complex Waring rank, and with \cref{thm_orbit}: in
  degree $6$ the icosahedral group has a two-dimensional invariant space, so a
  single orbit no longer suffices.  It is not a proof, since the search is
  nonconvex.
\end{remark}

%% ====================================================================
\section{One direction set for many operators}\label{sec_shared}

The implementation described in the main paper schedules an operator by
expanding it into monomials and scheduling each monomial separately.  The
following elementary observations quantify what that costs and identify the
right object to optimize.

\begin{definition}[Shared schedule]\label{def_shared}
  Let $\mathcal{P}_{1},\ldots,\mathcal{P}_{K}$ have symbols
  $\sigma_{1},\ldots,\sigma_{K} \in \K[\bz]_{p}$ and put
  $U := \spanop_{\K}\{\sigma_{1},\ldots,\sigma_{K}\}$.  A \emph{shared schedule}
  is a finite set $X = \{\bv_{1},\ldots,\bv_{R}\}$ together with coefficient
  vectors $c^{(i)} \in \K^{R}$ such that $(c^{(i)}_{r},\bv_{r})_{r}$ is a
  schedule for $\mathcal{P}_{i}$, for each $i$.  Write $R_{\K}(U)$ for the least
  such $R$.
\end{definition}

\begin{proposition}[Shared rank]\label{prop_shared}
  With $W_{X} := \spanop_{\K}\{(\bv\cdot\bz)^{p} : \bv \in X\}$, a shared
  schedule on $X$ exists if and only if $U \subseteq W_{X}$.  Moreover
  \begin{equation}\label{eq_shared_bounds}
    \max\Bigl(\dim_{\K} U,\ \max_{i} R_{\K}(\sigma_{i})\Bigr)
    \;\le\; R_{\K}(U)
    \;\le\; \min\Bigl(\sum_{i} R_{\K}(\sigma_{i}),\ D_{p,N}\Bigr).
  \end{equation}
\end{proposition}

\begin{proof}
  By \cref{prop_bridge} a schedule for $\mathcal{P}_{i}$ on $X$ exists exactly
  when $\sigma_{i} \in W_{X}$, and this for all $i$ is $U \subseteq W_{X}$ since
  $W_{X}$ is a subspace.  For the lower bounds, $\dim W_{X} \le \abs{X}$ gives
  $\dim U \le R_{\K}(U)$, and $U \subseteq W_{X}$ contains each $\sigma_{i}$ so
  $\abs{X} \ge R_{\K}(\sigma_{i})$.  For the upper bounds, the union of
  individually minimal sets is admissible, and any $X$ of size $D_{p,N}$ whose
  power images are independent has $W_{X} = \K[\bz]_{p} \supseteq U$.
\end{proof}

\begin{corollary}[The full derivative tensor]\label{cor_fulltensor}
  Taking $U = \K[\bz]_{p}$, that is, requiring every partial derivative of order
  $p$ in the $N$ active variables, gives
  \begin{equation}\label{eq_fulltensor}
    R_{\K}(\K[\bz]_{p}) \;=\; D_{p,N} \;=\; \binom{p+N-1}{N-1},
  \end{equation}
  attained by any $X$ of that cardinality whose power images are linearly
  independent, hence by a generic $X$.  No optimization is involved and no
  smaller set can work.
\end{corollary}

\begin{proof}
  Both bounds of \eqref{eq_shared_bounds} equal $D_{p,N}$.  Independence of
  $\{(\bv_{r}\cdot\bz)^{p}\}$ for generic $\bv_{r}$ holds because the Veronese
  variety is nondegenerate, so no nonzero linear functional on $\K[\bz]_{p}$
  annihilates all $p$th powers.
\end{proof}

\cref{cor_fulltensor} is the sharpest available answer to the efficiency
question when several derivatives are wanted at once: the correct cost is the
dimension of what is being asked for, not the sum of the individual ranks.  The
gap between the two is large already in modest cases; \cref{tab_numbers}
records it.

\begin{remark}[Conditioning is a separate objective]\label{rmk_cond}
  Minimal length does not imply minimal cost.  Let $M_{X}$ be the matrix of the
  map $c \mapsto \sum_{r}c_{r}(\bv_{r}\cdot\bz)^{p}$ in the monomial basis; the
  coefficients of a schedule are obtained by solving with $M_{X}$, so the
  relative error of the assembled derivative is governed by the condition number
  of $M_{X}$.  The Fourier grid of \cref{thm_grid} is a discrete Fourier matrix
  up to diagonal scaling and is extremely well conditioned, whereas generic
  minimal sets are not; \cref{tab_cond} measures both.  Orbit constructions are
  attractive precisely because they achieve short length and equal weights
  simultaneously, the icosahedral schedule of \cref{cor_tight} being minimal and
  perfectly balanced at once.  A realistic objective is therefore a weighted
  count of real and complex jet passes subject to a conditioning constraint,
  rather than the raw Waring rank.
\end{remark}

%% ====================================================================
\section{Numerical corroboration}\label{sec_numerics}

Every statement above that asserts a specific number has been checked
numerically by the script \texttt{verify\_theory\_directions.py} in
\texttt{experiments/tools}, which also
reproduces the tables of this section.  \cref{tab_numbers} compares, for
$\Delta^{m}$ and for the full derivative tensor, the length of the Fourier grid
of \cref{thm_grid}, the term-by-term total used by the current implementation,
the lower bound of \cref{thm_lower} or \cref{cor_fulltensor}, and the shortest
real schedule located by direct search.

\begin{table}[t]
\caption{Directions required, by construction.  The grid is
\cref{thm_grid} in the given coordinates; term-by-term sums the monomial ranks of
the main paper; the bound is \cref{thm_lower} for $\Delta^{m}$ and
\cref{cor_fulltensor} for the full tensor; the last column is the shortest real
schedule found by numerical search over nodes, with weights eliminated by
orthogonal projection.}
\label{tab_numbers}
\centering
\begin{tabular}{llrrrrl}
\toprule
target & $N$ & grid & term-by-term & bound & search & weights \\
\midrule
$\Delta^{1}$ & $3$ & $9$   & $3$   & $3$  & $3$  & positive \\
$\Delta^{2}$ & $3$ & $25$  & $12$  & $6$  & $6$  & positive, equal \\
$\Delta^{3}$ & $3$ & $49$  & $42$  & $10$ & $11$ & positive \\
$\Delta^{1}$ & $4$ & $27$  & $4$   & $4$  & $4$  & positive \\
$\Delta^{2}$ & $4$ & $125$ & $22$  & $10$ & $11$ & positive \\
\midrule
all order $4$ & $3$ & --- & $54$  & $15$ & $15$ & --- \\
all order $6$ & $3$ & --- & $192$ & $28$ & $28$ & --- \\
all order $4$ & $4$ & --- & $150$ & $35$ & $35$ & --- \\
all order $6$ & $4$ & --- & $776$ & $84$ & $84$ & --- \\
\bottomrule
\end{tabular}
\end{table}

Three readings of \cref{tab_numbers} matter.  First, the universal grid is a
poor choice for a structured operator: at $\Delta^{2}$ on $\R^{3}$ it needs $25$
directions where the term-by-term expansion needs $12$, so \cref{thm_grid} is a
fallback for symbols about which nothing is known, not a replacement for the
monomial schedule.  Second, term-by-term is nonetheless not optimal once the
operator couples coordinates: the same $\Delta^{2}$ needs only the six real
icosahedral axes, which \cref{cor_tight} certifies as minimal, and $\Delta^{3}$
on $\R^{3}$ drops from $42$ to $11$.  The saving grows with $m$ because the
term-by-term total grows like the number of monomials times their individual
ranks while the bound grows only like $D_{m,N}$.  Third, the largest gap appears
when many derivatives are wanted at once: the full order-six tensor in three
variables costs $192$ directions term by term against the exact optimum $28$ of
\cref{cor_fulltensor}, which requires no search at all.

\cref{tab_cond} quantifies \cref{rmk_cond}.  The grid buys its extra length back
in stability, which is the reason it should not be discarded even where shorter
schedules exist.

\begin{table}[t]
\caption{Length against conditioning for a symbol of degree $p$ in $N$
variables: a generic minimal set attaining \cref{cor_fulltensor} versus the
Fourier grid of \cref{thm_grid}.}
\label{tab_cond}
\centering
\begin{tabular}{rrrrrr}
\toprule
$p$ & $N$ & generic minimal & $\kappa(M_{X})$ & grid & $\kappa(M_{X})$ \\
\midrule
$4$ & $2$ & $5$  & $2.0\times10^{2}$ & $5$  & $6.0\times10^{0}$ \\
$6$ & $2$ & $7$  & $1.2\times10^{2}$ & $7$  & $2.0\times10^{1}$ \\
$4$ & $3$ & $15$ & $4.9\times10^{2}$ & $25$ & $1.2\times10^{1}$ \\
$6$ & $3$ & $28$ & $2.2\times10^{4}$ & $49$ & $9.0\times10^{1}$ \\
\bottomrule
\end{tabular}
\end{table}

%% ====================================================================
\section{Summary}\label{sec_summary}

The results assemble into a selection rule for building a schedule, ordered by
how much structure the symbol has.  If nothing is known about $\sigma$ beyond
its coordinate degrees, \cref{thm_grid} applies unconditionally, needs no solve,
and is well conditioned, at length $\prod_{j\ge1}m_{j}$.  If $\sigma$ is
projectively equivalent to a monomial, running the same construction in the
coordinates that make it a monomial is minimal, by \cref{rmk_tight}; this covers
the polyharmonic operator in two variables.  If $\sigma$ is a power of a
nondegenerate quadratic form, \cref{thm_cubature} replaces the design problem by
spherical cubature, giving real nodes, positive weights and, in the smallest
cases, provable minimality.  If $\sigma$ is invariant under any finite group,
\cref{thm_orbit} reduces the moment conditions to the dimension of the relevant
isotypic component.  And if several operators or a whole derivative tensor are
wanted, \cref{cor_fulltensor} gives an exactly optimal shared set of length equal
to the dimension of the span, with no search.

What is not settled is the general case.  For a symbol with no symmetry and no
quadratic structure in $N \ge 3$ active variables, no constructive procedure is
known that attains $R_{\C}(\sigma)$; computing the rank is already NP-hard, and
numerical decomposition can fail because the border rank may be strictly
smaller.  \cref{cor_price} makes the stakes concrete: the universal grid is off
by a factor tending to $N!$, so the gap is worth closing wherever structure
permits, and the four structured cases above are where it does.

\end{document}
